\documentclass[12pt]{article}
\usepackage[utf8]{inputenc}
\usepackage[T1]{fontenc}
\usepackage{hyperref}
\title{Wuding-Luquan Yi Language Resources for K-12 Curriculum Development}
\author{Compiled by Claude Sonnet 5.5 (claude-sonnet-5-5) for Andr\'as Kornai}
\date{}
\begin{document}
\maketitle

\section{Introduction}
Wuding-Luquan Yi (ISO 639-3 ywq, Glottocode wudi1238) is a Sino-Tibetan (Ngwi-Burmese) language of Yunnan, China. It is also known as Nasu (Wulu). The search found very little material that is usable in schools. This report lists only what the search returned. Sources disagree on speaker numbers: one gives about 10,000, while Wikipedia gives about 1.0 million for the wider Nasu grouping.

\section{Primary Educational Resources}
The search found no primers, textbooks or school programmes for this variety. Wikipedia notes that Nasu is written in the Pollard script, unlike most written Yi languages, but I found no classroom materials in it.

\section{Linguistic and Reference Materials}
\begin{itemize}
\item Wikipedia, Nasu language (general overview and classification): \url{https://en.wikipedia.org/wiki/Nasu_language}
\item Grambank entry for Wuding-Luquan Yi, which cites a grammar, \emph{Nasu yiyu yufa yanjiu} (Pu 2017): \url{https://grambank.clld.org/languages/wudi1238}
\item University of Hawai\textquoteleft i ScholarSpace archive item, a wordlist recording of Naisu made in Wuding County: \url{https://scholarspace.manoa.hawaii.edu/items/cf83ef6d-ec48-4ce8-bd00-d4e433c6eed8/full}
\end{itemize}

\section{Religious and Community Resources}
Global Recordings Network lists audio and audio-visual Bible-based lessons for this language, such as ``Good News'' and ``Words of Life'': \url{https://globalrecordings.net/en/language/ywq}

\section{Digital and Audio Resources}
The Global Recordings Network page above is the only audio source found. The same page is mirrored at \url{https://glorecnet-website.k.globalrecordings.net/en/language/ywq}.

\section{Curriculum Development Context}
No evidence was found that the language is taught in schools. A curriculum would have to be built from new work with community speakers, using the Pollard script or another orthography the community chooses. The 2017 grammar and the archived wordlist could be starting points for that work.
\end{document}